% !TeX root = ../../Book3.tex
\chapter{二连通区域与极值长度}
\label{b3:ch:08}
\BookChapterNumber{b3:ch:08}
\begin{chapterabstract}
圆环的半径比与矩形的宽高比都是共形形状参数，但仅凭区域的面积或外观很难判断它们。极值长度同时选择曲线族和非负密度，用最短长度的平方与面积之比测量通道的阻隔。它给出圆环与标记四边形的精确参数，并将本质包含、串联、并联和光滑拟共形变换转化为可计算的比较。圆环、去心盘与去心平面的比较还会显示，孔洞的尺度与端部类型怎样影响共形结构。
\end{chapterabstract}

\begin{learninggoals}
\begin{itemize}
\item\label{b3:item:08-goal-1}能在球面中辨认二连通域的两个补集分支，区分有限圆环、去心盘与去心平面，说明无限模量为何还需附带端部类型。
\item\label{b3:item:08-goal-2}能从密度定义证明极值长度的不变性和包含原则，核对长度与面积各自的换元，并辨认零面积修改对指定曲线的影响。
\item\label{b3:item:08-goal-3}能为圆环和矩形选择达到下界的密度，再用一族简单曲线证明上界，明确径向、绕行与指定对边的方向。
\item\label{b3:item:08-goal-4}能组合通道、选择试探势并比较能量，在光滑拟共形映射下给出模量界，区分必要下界、达到下界的构造和精确值。
\end{itemize}
\end{learninggoals}

两个圆盘无论大小都共形等价，两个圆环则未必。将圆环整体缩放，可以同时改变两个半径，却不能改变它们的比值；若允许一般全纯变换，这个比值是否仍然保留，就需要一个不依赖具体坐标的判据。上一章的矩形映射也留下同样的问题：四个直角不决定宽高比，而这个比值如何从原区域直接测出，不能总靠解出一条显式映射。

\ProseParagraphBreak

本章的方法是先指定允许的曲线，再分配密度。对于连接两端的曲线，我们关心的是所有路径中最省的长度；密度的整体倍数会同时改变长度和面积，因此取最短长度平方除以面积。固定面积预算后，把密度集中在每条路径必须经过的窄处，会使这个比值增大；若有许多并列路线，最短者又会限制这种分配。曲线族与密度承担不同的选择，先取下确界、再取上确界的次序不能交换。

\ProseParagraphBreak

\secref{b3:sec:0801}先辨认二连通模型，\secref{b3:sec:0802}引入环域模量及其比较。环域比较先陈述结论及几何例子，相关证明集中安排在\secref{b3:sec:0803}定义极值长度并算出圆环公式之后。这样先建立密度换元、Cauchy--Schwarz 与非负积分所给的工具，再证明参数唯一性与比较结论。\secref{b3:sec:0804}再将矩形计算传到标记四边形，并研究通道与光滑畸变。第七章的边界对应在传递四边形标记时使用；对一般粗糙环域，则应按内部端部及绕行曲线说明对应关系，不能预先假定普通边界逐点同胚。

\ProseParagraphBreak

单点补分支对应去心盘或去心平面，这些退化模型将帮助辨认无穷模量的几何含义。密度始终取明确的 Borel 代表，因为二维面积看不见的一条线，仍可能是曲线族中的一条完整路径。末节的拟共形估计只使用保向 \(C^1\) 微分同胚；更弱正则性下的曲线例外及延拓问题留待附录B。

% !TeX root = ../../Book3.tex
\section{二连通区域}
\label{b3:sec:0801}

单连通真平面区域只有圆盘这一内部共形模型。二连通区域开始出现一个不能靠归一化消去的实参数。先区分标准模型与退化类型，再以曲线族给这个参数一个可计算的定义。

\subsection{环域}
\label{b3:subsec:080101}

\begin{definition}\label{b3:def:08-doubly-connected}
扩充平面中区域 \(U\) 称为二连通，若 \(\widehat{\mathbb C\kern0pt}\setminus U\) 恰有两个非空连通分支。每个分支是一个紧连通集，允许退化为单点。
\end{definition}
本章“环域”在广义讨论中指这种区域，标准圆环则另加“圆”字。二连通并不要求两条光滑边界；裂缝、线段或复杂紧连通集都可作为补集分支。拓扑上的环与共形标准圆环也不同：拓扑等价不控制两端之间的解析距离。

\ProseParagraphBreak
\(\mathbb C\setminus\{0\}\) 在球面中删去了0和无穷远，是二连通；单位去心盘的两个补集分支是0与闭外部。普通环 \(r<|z|<R\) 的两个补集分支都是非退化的。要在球面中计算补集，才能给无界区域作一致分类。

\subsection{标准环域}
\label{b3:subsec:080102}

标准圆环记为 \(A(r,R)=\{r<|z|<R\}\)，其中 \(0<r<R<\infty\)。缩放可令外半径为1。映射 \(z\mapsto rR/z\) 交换两端，旋转保持每一端。去心盘 \(\mathbb D^*\) 与去心平面 \(\mathbb C^*\) 是两个退化模型。

\ProseParagraphBreak
取对数后，圆环的万有覆盖是竖直带 \(\log r<\operatorname{Re}w<\log R\)，甲板平移为 \(w\mapsto w+2\pi\mathrm{i}\)。这是直接由指数的周期验证的模型，不需要先使用一般万有覆盖定理。带宽与平移长度的比将成为模量。去心盘对应半平面 \(\operatorname{Re}w<0\)，去心平面对应整个 \(w\) 平面。

\subsection{分类定理}
\label{b3:subsec:080103}

\begin{theorem}[深层二连通分类]\label{b3:thm:08-annulus-classification}
二连通区域 \(U\) 的两个补集分支若都非单点，则 \(U\) 双全纯等价于唯一的 \(A(r,1)\)、\(0<r<1\)。若恰有一个分支为点，则等价于 \(\mathbb D^*\)；若两个都为点，则等价于 \(\mathbb C^*\)。
\end{theorem}
\begin{proofsketch}
任意边界下的存在部分是本章明确采用的深层输入。可先建立万有覆盖与统一化，将二连通域写成单连通模型对无限循环甲板群的商；圆盘模型中的生成元有抛物或双曲类型，其商分别是去心盘与有限圆环，平面模型的平移商为去心平面。还要证明点状补集分支与退化端一致，不能仅凭基本群为整数群区分三者。此路线在第11--13章解释；一般二连通映射可参阅 Ahlfors《Conformal Invariants》\cite{Ahlfors2010Conformal}关于环域与极值长度的论述。标准参数的唯一性在下文由完整的极值长度计算证明，未作为未经说明的假定。
\end{proofsketch}
分类的存在性并不意味着给出显式映射。对于两条不同心圆之间的区域，分式变换可先把两圆同心化；对于一般裂缝环域，常需使用数值或势论方法。

\subsection{共形不变量}
\label{b3:subsec:080104}

有限圆环的半径比是一个真正的共形不变量。同样面积或同样欧氏厚度的两个圆环可有不同半径比，因此不共形等价。另一方面，\(A(r,R)\) 与 \(A(cr,cR)\) 对任意 \(c>0\) 等价，绝对大小并非不变量。

\ProseParagraphBreak
退化类型还需单独记录。去心盘与去心平面都具有下文意义的无限模量，却不等价：若有 \(\mathbb C^*\rightarrow\mathbb D^*\) 双全纯映射，作为有界函数它在0可去，成为有界整函数，必常值。无限模量不能独自区分这两个端部结构。

\ProseParagraphBreak
“二连通”在这里以球面补集的连通分支数定义。对于 \(A(r,R)\)，两个补分支分别是内闭盘和外闭盘加无穷远，都非单点；去心盘的内分支退化为一个点，外分支仍非单点；去心平面的两分支都为点。圆盘本身在球面中只有一个补分支，虽在某幅平面图里可以同时看到“外面”和“无穷远”，它们属于同一个分支，不能重复计数。

\ProseParagraphBreak
两端的拓扑资料还可从穷竭看出。去心盘中离穿孔越来越近的一列小圆与离外圆越来越近的一列大圆代表两个不同端；前者在球面闭包收缩到一个点，后者的极限集是整个外边界。双全纯映射保持内部端的趋近结构，却不必把任意粗糙普通边界逐点对应。这是正文用端部曲线族定义模量、而没有在一般分类里提前假定边界同胚的原因。

% !TeX root = ../../Book3.tex
\section{环域模量}
\label{b3:sec:0802}

模量量化两个边界分支之间的相对分离。它可以通过标准圆环定义，也可以不用求映射，直接由曲线族的长度与面积极值计算。后一种定义使比较和退化都可证明。

\subsection{环域模量}
\label{b3:subsec:080201}

\begin{definition}\label{b3:def:08-annulus-modulus}
若二连通区域 \(U\) 等价于 \(A(r,R)\)，定义\term[huanyu-moliang]{环域模量}[modulus of an annulus]为
\[
 \operatorname{mod}U=\frac1{2\pi}\log\frac Rr.
\]
对去心盘与去心平面记模量为 \(+\infty\)，并另记其退化类型。
\end{definition}
\symidx{\(\operatorname{mod}U\)：环域模量}
宽圆环的模量大，薄圆环的模量小；这与径向连接两边界的曲线“更长”的直觉一致。\(1/(2\pi)\) 来自绕行周长的归一化，不可在后面的矩形计算中随意删去。

\subsection{共形不变性}
\label{b3:subsec:080202}

\begin{proposition}\label{b3:prop:08-modulus-invariance}
两个有限圆环双全纯等价，当且仅当它们半径比相同。
\end{proposition}
\begin{proof}
必要性用本章第三节将完整证明的两项事实：极值长度在双全纯变换下不变，且连接圆环两端的曲线族的极值长度为 \((2\pi)^{-1}\log(R/r)\)。边界端由区域内部的趋近方式刻画，所以双全纯映射将此族对应到目标连接族，允许交换两端。充分性由缩放或倒数直接构造。
\end{proof}
这里证明只向前引用本章独立计算，不借模量不变性来证明极值长度不变性，因而没有循环。分类定理中的“唯一半径”也从此得到。

\subsection{单调性}
\label{b3:subsec:080203}

设 \(U\subset V\) 均为非退化二连通域，并且 \(U\) 中绕孔一周的曲线在 \(V\) 仍绕孔一周，称该包含为本质包含。由绕行曲线族的包含与下一节的倒数公式，有
\[
 \operatorname{mod}U\le\operatorname{mod}V.
\]
本质性必不可少。一个很小的环可以嵌在大环的局部圆盘内，它的孔在大环中可以被填掉；此时两者的模量没有由集合包含产生的固定关系。

\ProseParagraphBreak
具体地，\(A(r_2,R_2)\subset A(r_1,R_1)\) 若 \(r_1\le r_2<R_2\le R_1\)，中心孔相同，确为本质包含，模量不等式就是半径比的比较。对于非圆边界，只要能把区域夹在两个本质同心圆环之间，也可取得无需显式映射的上下界。

\begin{proposition}[串联环域的模量比较]\label{b3:prop:08-series-annuli}
设若干互不相交的二连通子域 \(U_1,\ldots,U_N\) 在环域 \(U\) 的两端之间依次分隔，且每条连接 \(U\) 两端的曲线都包含依次穿过各 \(U_j\) 的连接子曲线。则
\[
 \operatorname{mod}U\ge\sum_{j=1}^N\operatorname{mod}U_j.
\]
\end{proposition}
\begin{proof}
记 \(m_j=\operatorname{mod}U_j\)。由极值长度定义，在 \(U_j\) 上选密度 \(\rho_j\)，面积为 \(A_j>0\)，连接族的最短长度为 \(L_j>0\)，使 \(L_j^2/A_j\) 任意接近 \(m_j\)。把 \(\rho_j\) 乘 \(L_j/A_j\)，在其余区域置0，得到 \(U\) 上密度。每条总连接曲线长度至少为 \(\sum_j L_j^2/A_j\)，而总面积恰为这个和。因此总族极值长度至少为该和。逐一逼近各上确界便得所述式。零模量或极端情况先取有限近似，同样成立。
\end{proof}
同心圆环 \(A(r,s)\) 和 \(A(s,R)\) 在 \(A(r,R)\) 中串联，等号成为对数恒等式 \(\log(R/r)=\log(s/r)+\log(R/s)\)。中间圆周不属两个开环，但零面积的分隔曲线不改变所用长度积分。这项比较揭示细长瓶颈可以相加：只要每条总连接曲线必须穿过各段，局部阻碍便累积。若子环并列放置而曲线可以择一通过，串联假设失败，不能照搬求和。

\ProseParagraphBreak
可以给非本质包含一个数值反例。固定大环 \(V=A(1,3)\)，在点2附近放一个外半径为 \(\varepsilon<1/2\) 的小环 \(U=A_2(\varepsilon e^{-2\pi M},\varepsilon)\)，即把同心环平移到中心2。它完全包含于 \(V\)，模量恰为任意指定 \(M>0\)。当 \(M>\log3/(2\pi)\) 时，模量比大环还大；把半径比取接近1，又可任意小。小环的孔位于 \(V\) 内部，绕小孔的圈在 \(V\) 可收缩，所以该包含不是本质包含。由集合大小直接判断模量大小，会在此例失败。

\ProseParagraphBreak
串联比较也可使用极值密度的归一化解释。每段给一个最短连接长度1、面积接近 \(1/m_j\) 的密度。将其乘 \(m_j\)，长度贡献成为 \(m_j\)，面积贡献也成为 \(m_j\)。总路径必须付出各段长度之和，而总面积只按各段相加，因此总商至少为总模量。这相当于电阻串联，但证明并不依赖物理比喻；不相交支集与必须逐段穿过的条件已经给出全部数学依据。

\subsection{退化}
\label{b3:subsec:080204}

当 \(r\downarrow0\) 而 \(R\) 固定时，有限圆环模量趋于无穷，内端收缩为去心盘的点端。若 \(R/r\downarrow1\)，模量趋于0，径向厚度相对周长消失。两种极限分别对应长圆柱与薄圆柱；它们不应混称为同一种“退化”。

\ProseParagraphBreak
若同时令 \(r\downarrow0,R\uparrow\infty\)，局部区域穷竭为去心平面，模量也趋于无穷。与固定 \(R\) 的退化相比，第二端也走向无穷远，极限类型因此不同。只保存一个实数参数而不保存区域的观察点和端部数据，会在退化极限中丢失信息。

% !TeX root = ../../Book3.tex
\section{极值长度}
\label{b3:sec:0803}

极值长度以所有非负密度为竞争者，比的是最短曲线长度的平方与面积。缩放密度不改变比值，因此它是一个没有欧氏单位的量。

\subsection{曲线族}
\label{b3:subsec:080301}

曲线族 \(\Gamma\) 是区域内局部可求长曲线的集合。本章圆环的径向族包含所有连接两端的曲线，允许以边界端作为极限端点；绕行族包含绕中心指数为1或 \(-1\) 的闭曲线。矩形族则连接指定两条相对边。

\ProseParagraphBreak
同一个区域配不同曲线族会给不同极值量。径向族和绕行族将得到互为倒数的数值。曲线参数化不重要，长度积分使用弧长；重复往返只增加非负长度，不会降低连接端点所需的最短量。允许曲线不光滑并不妨碍定义，只需按局部可求长性理解积分。

\subsection{可容许度量}
\label{b3:subsec:080302}

取 Borel 可测密度 \(\rho\ge0\)，要求
\[
 0<A_\rho=\int_U\rho^2\,\mathrm{d}A<\infty.
\]
定义 \(L_\rho(\gamma)=\int_\gamma\rho\,\mathrm{d}s\)，以及 \(L_\rho(\Gamma)=\inf_{\gamma\in\Gamma}L_\rho(\gamma)\)。零密度区域允许存在，个别曲线长度也允许无穷。极值问题只比较有限正面积的密度，不把 \(0/0\) 或 \(\infty/\infty\) 当作竞争值。

\ProseParagraphBreak
与“令每条曲线长度至少1，再最小化面积”的表述相比，本章不先使用可容许一词作为额外条件；若 \(L_\rho(\Gamma)>0\)，把密度除以这个数即可转成那种等价形式。若族里有任意短曲线使下确界为0，该密度的贡献就是0。

\ProseParagraphBreak
极值密度的整体倍数不影响商。将 \(\rho\) 乘 \(c>0\)，最短长度乘 \(c\)，面积乘 \(c^2\)，所以 \(L^2/A\) 不变。面积可以归一化为1，也可把最短长度归一化为1。在后一归一化中
\[
 \frac1{\lambda(\Gamma)}=\inf\left\{\int\rho^2\,\mathrm{d}A:
       \int_\gamma\rho\,|\mathrm{d}z|\ge1\quad(\gamma\in\Gamma)\right\},
\]
只在 \(0<\lambda<\infty\) 时直接写此互倒公式。它来自对每个密度作尺度调整，而不额外假设上确界达到。这个面积下确界也常被称为曲线族的模；不同文献的“模”可能指极值长度或其倒数，使用时必须先核对定义。

\ProseParagraphBreak
非负 Borel 密度可以在某些点为0或无穷，只要求总面积有限且非零。对零面积集合任意改变密度未必不影响每一条曲线的长度，例如一条恰落在该集合中的曲线；因此若把密度视为 \(L^2\) 等价类，就还要加入曲线族例外的处理。本章初等计算使用明确的密度代表，长度不等式按实际曲线验证，不把平面几乎处处相等自动转成所有路径的相等。

\subsection{极值长度}
\label{b3:subsec:080303}

\begin{definition}\label{b3:def:08-extremal-length}
\term[jizhi-changdu]{极值长度}[extremal length]定义为
\[
 \lambda_U(\Gamma)=\sup_\rho\frac{L_\rho(\Gamma)^2}{A_\rho}.
\]
\end{definition}
\begin{theorem}\label{b3:thm:08-extremal-conformal}
若 \(f:U\rightarrow V\) 双全纯，则 \(\lambda_U(\Gamma)=\lambda_V(f\Gamma)\)。若 \(\Gamma_1\subset\Gamma_2\)，则 \(\lambda(\Gamma_1)\ge\lambda(\Gamma_2)\)。
\end{theorem}
\begin{proof}
给目标密度 \(\sigma\)，拉回 \(\rho(z)=\sigma(f(z))|f'(z)|\)。链式法则给每条曲线长度相等，平面换元的 Jacobi 行列式为 \(|f'|^2\)，所以面积相等。逆映射把任意源密度也这样对应到目标密度，取上确界得相等。曲线族包含时，下确界在较小集合上较大，对每个密度比较后取上确界即得第二式。
\end{proof}
乘密度一个正数使长度与面积分别乘该数及其平方，故比值保持不变。真正重要的是密度的分布而非总尺度。

\begin{proposition}[圆环的 Dirichlet 能量]\label{b3:prop:08-annulus-capacity}
设 \(A=A(r,R)\)，\(u\) 为闭圆环上的 \(C^1\) 实函数，在内圆取0、外圆取1。则
\[
 \int_A|\nabla u|^2\,\mathrm{d}A\ge\frac{2\pi}{\log(R/r)}.
\]
等号由 \(u_0(z)=\log(|z|/r)/\log(R/r)\) 达到，并且在这类函数中唯一。
\end{proposition}
\begin{proof}
记 \(L=\log(R/r)\)。沿每条径向线，\(\int_r^R\partial_su\,\mathrm{d}s=1\)。Cauchy--Schwarz 给
\[
 1\le\left(\int_r^R |\partial_su|^2s\,\mathrm{d}s\right)
       \left(\int_r^R\frac{\mathrm{d}s}{s}\right).
\]
对角度积分并舍去非负角方向能量，得到下界 \(2\pi/L\)。\(u_0\) 的梯度为径向 \(1/(Ls)\)，直接积分给等号。等号必须使角导数为0，且各径向 Cauchy--Schwarz 等号使 \(\partial_su\) 为 \(c/s\)；边界值迫使 \(c=1/L\)，故只能为 \(u_0\)。
\end{proof}
这个最小能量称为圆环的容量，它恰是径向连接族极值长度的倒数，也恰是绕行族的极值长度。若把圆环切开并用 \(w=\log z\)，\(u_0\) 变成宽度为 \(L\) 的平直圆柱上的线性势。几何比较、曲线长度和能量三个观点在此给同一数值。

\ProseParagraphBreak
平面共形换元使实函数的 Dirichlet 能量不变：梯度模平方乘 \(|f'|^{-2}\)，面积元乘 \(|f'|^2\)，两项抵消。于是对光滑边界环域可通过环域映射把最小能量带过去；边界值的对应须先有合适的边界延拓。对于任意粗糙边界，能量边界条件需要另行定义，不能直接写“在每个边界点取0或1”并假定经典解存在。后面的势论将系统研究这种问题。

\subsection{环域计算}
\label{b3:subsec:080304}

\begin{theorem}\label{b3:thm:08-annulus-extremal}
在 \(A(r,R)\) 中，连接两端的族 \(\Gamma_{\mathrm r}\) 与绕行族 \(\Gamma_{\mathrm c}\) 满足
\[
 \lambda(\Gamma_{\mathrm r})=\frac{\log(R/r)}{2\pi},\qquad
 \lambda(\Gamma_{\mathrm c})=\frac{2\pi}{\log(R/r)}.
\]
\end{theorem}
\begin{proof}
取 \(\rho=1/|z|\)，面积为 \(2\pi\log(R/r)\)。任意径向连接曲线的长度至少 \(\log(R/r)\)，由沿曲线积分 \(|\mathrm{d}\log|z||\le|\mathrm{d}z|/|z|\) 得到；径向线达到此值，所以给第一式下界。任意绕行曲线长度至少 \(2\pi\)，由辐角总变化得到；圆周达到此值，给第二式下界。

任取密度 \(\sigma\)。径向线的长度 \(\ell(\theta)=\int_r^R\sigma(se^{\mathrm{i}\theta})\,\mathrm{d}s\) 由 Cauchy--Schwarz 满足
\[
 \ell(\theta)^2\le\log(R/r)\int_r^R\sigma(se^{\mathrm{i}\theta})^2s\,\mathrm{d}s.
\]
对角度平均，存在径向线使长度平方不超过 \(A_\sigma\log(R/r)/(2\pi)\)，得到第一式上界。对圆周长度 \(h(s)=\int_0^{2\pi}\sigma(se^{\mathrm{i}\theta})s\,\mathrm{d}\theta\)，有 \(h(s)^2\le2\pi s^2\int_0^{2\pi}\sigma^2\,\mathrm{d}\theta\)。再以 \(\mathrm{d}s/s\) 积分，存在半径使 \(h(s)^2\le2\pi A_\sigma/\log(R/r)\)，得到第二式上界。可测密度由 Tonelli 定理保证这些积分对几乎所有参数成立；取下确界不受零测参数集影响。
\end{proof}
极值密度的几何是把对数圆柱展开后的常密度。两族一纵一横，长度平方除面积互为倒数。

\begin{proposition}[试探势给极值长度下界]\label{b3:prop:08-test-potential}
设有界平面通道 \(Q\) 的两个指定边界段为 \(E_0,E_1\)，\(u\) 在其闭包附近 \(C^1\)，在两段分别取0、1，且 \(0<E=\int_Q|\nabla u|^2\,\mathrm{d}A<\infty\)。连接这两段的曲线族满足 \(\lambda(\Gamma)\ge1/E\)。
\end{proposition}
\begin{proof}
取密度 \(\rho=|\nabla u|\)。任一连接曲线由链式法则给 \(1=|u(\gamma(1))-u(\gamma(0))|\le\int_\gamma\rho\,|\mathrm{d}z|\)，所以最短长度至少1，总面积为 \(E\)，代入定义即得下界。
\end{proof}
圆环的对数势达到最佳下界，矩形的线性势也达到最佳下界。一般通道可先构造一个方便的 \(u\)，即使它不是调和函数，仍给严格的下界。越小的能量给越强的下界；寻找最小能量自然引向 Dirichlet 问题，但这个简单比较本身不需要先证明最小解存在。若势只具有弱导数，沿所有曲线的链式法则需再处理，不能把本命题的 \(C^1\) 条件无声地删除。

% !TeX root = ../../Book3.tex
\section{四边形与共形模}
\label{b3:sec:0804}

圆环用两个边界分支区分连接方向，单连通区域则可以用四个边界标记指定两对相对边。把这样一个区域化为矩形以后，极值长度就是宽高比。本节先证明这个参数如何由曲线族确定，再用密度构造估计变宽通道和相互隔开的并联通道。最后考察非共形的光滑变换：长度沿不同方向伸缩不一，但最大、最小伸长之比仍能控制共形模的变化。

\subsection{四边形}
\label{b3:subsec:080401}

本节的标记四边形由一个有界 Jordan 区域 \(Q\) 和四个互异边界点 \(p_0,p_1,p_2,p_3\) 组成，四点按正向，也就是逆时针次序排列。所谓四边形只表示存在四个标记，不要求边界由直线组成。为与\cref{b3:prop:07-elliptic-rectangle}的矩形保持一致，我们规定四点在标准矩形上依次为左下、右下、右上、左上角。于是沿正向边界的四条开弧分别是
\[
 (p_0,p_1):\text{下边},\quad
 (p_1,p_2):\text{右边},\quad
 (p_2,p_3):\text{上边},\quad
 (p_3,p_0):\text{左边}.
\]
这里 \((p_i,p_j)\) 表示指定方向的边界弧，端点不计入开弧。如果依次说“左、上、右、下”，那是顺时针的读法，不能同时当作上述顶点的正向编号。在一般区域中，左、右等词只是边界标记的名称，不表示这些弧在原平面的实际方位。

\ProseParagraphBreak
记 \(\Gamma_{\mathrm h}\) 为内部局部可求长、两极限端点分别落在左、右开边上的全部曲线；\(\Gamma_{\mathrm v}\) 同样连接下、上开边。参数曲线可以在接近边界时有无限总长度，只要求每条内部紧子弧可求长。四个标记及连接哪一对边都是问题的数据。把标记循环移一位会交换两族，移两位则只交换同一对边的两端。

\Needspace{6\baselineskip}
\begin{proposition}[保持标记的矩形映射]\label{b3:prop:08-marked-rectangle-existence}
设 \(Q\subseteq\mathbb C\) 为有界 Jordan 区域，\(p_0,p_1,p_2,p_3\in\partial Q\) 互异并按正向排列。则存在 \(a,b>0\) 及双全纯映射
\[
 F:Q\longrightarrow R_{a,b}
 \coloneqq\{\,x+\mathrm{i}y:0<x<a,\ 0<y<b\,\},
\]
使 \(F\) 延为闭区域间的同胚，且
\[
 F(p_0)=0,\quad F(p_1)=a,\quad
 F(p_2)=a+\mathrm{i}b,\quad F(p_3)=\mathrm{i}b.
\]
\end{proposition}
\begin{proof}
Riemann 映射定理先将 \(Q\) 映到圆盘。因为 \(Q\) 为 Jordan 区域，\cref{b3:thm:07-caratheodory-jordan}保证此映射有边界同胚。再用圆盘到上半平面的分式变换和上半平面自同构，将 \(p_0,p_1,p_3\) 的像依次归一化为 \(0,1,\infty\)。这些变换保持边界的循环方向，所得映射 \(\psi:Q\rightarrow\mathbb H\) 必满足
\[
 \psi(p_0)=0,\quad\psi(p_1)=1,\quad
 \psi(p_2)=\tau>1,\quad\psi(p_3)=\infty.
\]
不等式 \(\tau>1\) 来自第四点在 \(1\) 与 \(\infty\) 之间的边界次序，并非对原区域欧氏坐标的要求。

取 \(k=\tau^{-1/2}\in(0,1)\)、\(k'=\sqrt{1-k^2}\)。\cref{b3:prop:07-elliptic-rectangle}已证明，椭圆积分映射 \(F_k\) 把预顶点 \(0,1,k^{-2},\infty\) 依次映到
\[
 0,\quad K(k),\quad K(k)+\mathrm{i}K(k'),\quad
 \mathrm{i}K(k'),
\]
并在闭包上给出边界同胚。因此 \(F=F_k\circ\psi\) 满足结论，其中 \(a=K(k)\)、\(b=K(k')\)。再乘任意正数，可以同时改变两个边长而保持标记。
\end{proof}

上述构造保证矩形模型存在，还没有说明不同构造是否得到相同宽高比。这个唯一性将由矩形中两族曲线的精确极值长度得到。若边界不是 Jordan 曲线，四个普通边界点未必代表四个不同趋近方向；此时须改用四个循环有序的素端，并使用素端边界上的对应。本节先在 Jordan 情形中完成计算和证明。

\subsection{对偶曲线族与通道比较}
\label{b3:subsec:080402}

在矩形里，连接左右边必须付出至少一个宽度，连接上下边必须付出至少一个高度。常密度因而给出两个下界。上界不能只挑一条方便的曲线，而要说明对任意密度，总能从水平段或竖直段中找到足够短的曲线；面积积分正好能完成这种平均。

\Needspace{6\baselineskip}
\begin{theorem}[矩形的两个连接族]\label{b3:thm:08-rectangle-extremal}
设 \(a,b>0\)，\(R_{a,b}=(0,a)\times(0,b)\subseteq\mathbb C\)。令 \(\Gamma_{\mathrm h}\) 为 \(R_{a,b}\) 内局部可求长、两极限端点分别在左、右开边上的全部曲线，\(\Gamma_{\mathrm v}\) 为同样连接下、上开边的全部曲线。则
\[
 \lambda_{R_{a,b}}(\Gamma_{\mathrm h})=\frac ab,\quad
 \lambda_{R_{a,b}}(\Gamma_{\mathrm v})=\frac ba.
\]
常密度 \(\rho\equiv1\) 同时达到两个上确界。
\end{theorem}
\begin{proof}
对常密度，任意左右连接曲线的欧氏长度至少为端点横坐标之差 \(a\)，任意上下连接曲线的长度至少为纵坐标之差 \(b\)。对只在内部局部可求长的曲线，先在紧子弧上使用这个距离下界，再令两端趋近边界即可。水平段、竖直段分别达到这两个长度，故
\[
 L_1(\Gamma_{\mathrm h})=a,\quad
 L_1(\Gamma_{\mathrm v})=b,\quad A_1=ab.
\]
定义中的两个商分别为 \(a/b\)、\(b/a\)，这就给出下界。

反过来，给定非负 Borel 密度 \(\rho\)，设 \(0<A=\int_{R_{a,b}}\rho^2\,\mathrm dA<\infty\)。记水平段与竖直段的密度长度为
\[
 q_{\mathrm h}(y)=\int_0^a\rho(x,y)\,\mathrm dx,\quad
 q_{\mathrm v}(x)=\int_0^b\rho(x,y)\,\mathrm dy.
\]
沿每条线段使用 Cauchy--Schwarz，再由 Tonelli 定理交换非负积分，得到
\[
 \begin{aligned}
 \int_0^b q_{\mathrm h}(y)^2\,\mathrm dy
 &\le a\int_0^b\int_0^a\rho(x,y)^2\,\mathrm dx\,\mathrm dy=aA,\\
 \int_0^a q_{\mathrm v}(x)^2\,\mathrm dx
 &\le b\int_0^a\int_0^b\rho(x,y)^2\,\mathrm dy\,\mathrm dx=bA.
 \end{aligned}
\]
特别地，两类线段的长度分别对几乎处处的参数有限。每条水平段属于 \(\Gamma_{\mathrm h}\)，每条竖直段属于 \(\Gamma_{\mathrm v}\)，所以
\[
 b\cdot L_\rho(\Gamma_{\mathrm h})^2\le aA,\quad
 a\cdot L_\rho(\Gamma_{\mathrm v})^2\le bA.
\]
除以正面积 \(A\)，再对所有密度取上确界，便得到两个匹配的上界。个别线段的长度可能无穷，但它们只出现在一个参数零测集上，不影响上述积分论证，也不需要把密度修改为另一个几乎处处相等的代表。
\end{proof}

两个极值长度的乘积为 \(1\)，现在是由同一张矩形的两个完整计算得出的事实。对一般标记四边形，保持标记的矩形映射把两族分别送到这两个标准族；\cref{b3:thm:08-extremal-conformal}的共形不变性再把倒数关系传回原区域。其他两族曲线即使都位于同一区域，也不能因此被任意称为这种对偶族。

\ProseParagraphBreak
无需先求出矩形映射，也常能从通道宽度取得有效估计。变宽通道中，较窄截面会提高连接两端的代价，因此密度应在这些位置更大。取宽度的倒数，长度与面积就会出现同一个积分。

\Needspace{6\baselineskip}
\begin{proposition}[变宽通道的下界与夹逼]\label{b3:prop:08-variable-channel}
设 \(a,b\in\mathbb R\)、\(a<b\)，\(h:[a,b]\rightarrow(0,\infty)\) 连续，并令
\[
 Q=\{\,(x,y):a<x<b,\ 0<y<h(x)\,\},\quad
 h_-=\min_{[a,b]}h,\quad h_+=\max_{[a,b]}h.
\]
令 \(\Gamma\) 为 \(Q\) 内所有局部可求长、两极限端点分别落在左、右开边上的曲线族，其中左边为 \(\{a\}\times(0,h(a))\)，右边为 \(\{b\}\times(0,h(b))\)。则
\[
 \frac{b-a}{h_+}
 \le\int_a^b\frac{\mathrm dx}{h(x)}
 \le\lambda_Q(\Gamma)
 \le\frac{b-a}{h_-}.
\]
\end{proposition}
\begin{proof}
连续性及正性保证 \(0<h_-\le h_+<\infty\)。记 \(I=\int_a^b h(x)^{-1}\,\mathrm dx\)，取 \(\rho(x,y)=1/h(x)\)。它的面积为
\[
 A_\rho=\int_a^b\int_0^{h(x)}\frac{\mathrm dy\,\mathrm dx}{h(x)^2}=I.
\]
要估计所有连接曲线，令 \(H(x)=\int_a^x h(t)^{-1}\,\mathrm dt\)。在一条内部紧子弧的弧长参数 \(s\) 下，横坐标 \(x(s)\) 绝对连续且 \(|x'(s)|\le1\)。链式法则因此给出
\[
 \int_{s_0}^{s_1}\frac{\mathrm ds}{h(x(s))}
 \ge\int_{s_0}^{s_1}\frac{|x'(s)|}{h(x(s))}\,\mathrm ds
 \ge\bigl|H(x(s_1))-H(x(s_0))\bigr|.
\]
让子弧两端分别趋近左右边界，便有 \(L_\rho(\gamma)\ge I\)。这也允许曲线在横向往返，往返只会增加第一个积分。另一方面，水平线 \(y=h_-/2\) 完全位于 \(Q\)，其密度长度恰为 \(I\)。所以 \(L_\rho(\Gamma)=I\)，得到 \(\lambda_Q(\Gamma)\ge I^2/I=I\)。

再取矩形 \(R_-=(a,b)\times(0,h_-)\)、\(R_+=(a,b)\times(0,h_+)\)，记其左右连接族为 \(\Gamma_-\)、\(\Gamma_+\)。有
\[
 R_-\subseteq Q\subseteq R_+,\quad
 \Gamma_-\subseteq\Gamma\subseteq\Gamma_+.
\]
\cref{b3:prop:08-ambient-independence}保证，同一个曲线族放在较大的开集里计算，极值长度不变；在原区域外把密度置零即可实现反向比较。因而可以把三族都看作 \(R_+\) 内的族，再用曲线族包含原则及矩形计算，得到
\[
 \frac{b-a}{h_+}\le\lambda_Q(\Gamma)\le\frac{b-a}{h_-}.
\]
最后，由 \(h(x)\le h_+\) 直接积分得 \(I\ge(b-a)/h_+\)，合并即得陈述中的全部不等式。
\end{proof}

常宽时这些界同时成为等式。变宽时，我们已经找出所选密度的最短曲线及其精确商 \(I\)，但尚未证明这个密度使所有竞争者的商最大。因此一般只有积分下界；通道上下边的变化可能让另一种密度给出更大的商。

\ProseParagraphBreak
如果可以从若干相互隔开的通道中任选一条，允许的曲线族变大，极值长度会下降。对于真正不相交的通道，下降幅度还有一个精确公式。

\Needspace{6\baselineskip}
\begin{proposition}[不交通道的并联公式]\label{b3:prop:08-parallel-channels}
设正整数 \(N\ge1\)，\(Q_1,\ldots,Q_N\subseteq\mathbb C\) 为两两不交的非空开集。对每个 \(j\)，\(\Gamma_j\) 是完全位于 \(Q_j\) 内的局部可求长曲线族，满足
\[
 0<\lambda_j\coloneqq\lambda_{Q_j}(\Gamma_j)<\infty.
\]
令 \(U=\bigcup_{j=1}^NQ_j\)、\(\Gamma=\bigcup_{j=1}^N\Gamma_j\)。则
\[
 \frac1{\lambda_U(\Gamma)}=\sum_{j=1}^N\frac1{\lambda_j}.
\]
\end{proposition}
\begin{proof}
给定 \(U\) 上密度 \(\rho\)，设 \(0<A_\rho<\infty\)。记它在各通道上的面积和最短长度为
\[
 A_j=\int_{Q_j}\rho^2\,\mathrm dA,\quad
 L_j=\inf_{\gamma\in\Gamma_j}\int_\gamma\rho\,\mathrm ds.
\]
不相交性给 \(A_\rho=\sum_jA_j\)，而有限并集的下确界给 \(L_\rho(\Gamma)=\min_jL_j\)。若某个 \(A_j=0\)，由 \(\lambda_j<\infty\) 及\cref{b3:lem:08-zero-area-density}可知 \(L_j=0\)，于是总最短长度为 \(0\)，所需上界成立。若所有 \(A_j>0\)，则定义给出 \(L_j^2\le\lambda_j A_j\)。令 \(L=\min_jL_j\)，便有
\[
 L^2\sum_{j=1}^N\frac1{\lambda_j}
 \le\sum_{j=1}^N\frac{L_j^2}{\lambda_j}
 \le\sum_{j=1}^NA_j=A_\rho.
\]
所以每个竞争商都不超过 \((\sum_j\lambda_j^{-1})^{-1}\)，取上确界得到一个方向。

反向不需要假定各通道存在极值密度。给定 \(\varepsilon>0\)，由 \(0<\lambda_j<\infty\) 的上确界定义，可在每个通道选择密度，其商为正且足够接近 \(\lambda_j\)，使这个商的倒数小于 \(\lambda_j^{-1}+\varepsilon/N\)。此密度的最短长度有限且为正，将它除以自己的最短长度，得到密度 \(\rho_j\)，满足
\[
 L_{\rho_j}(\Gamma_j)=1,\quad
 0<A_{\rho_j}<\frac1{\lambda_j}+\frac\varepsilon N.
\]
在每个 \(Q_j\) 放置 \(\rho_j\)，就得到 \(U\) 上的一个 Borel 密度。整体最短长度为 \(1\)，总面积小于 \(\sum_j\lambda_j^{-1}+\varepsilon\)。因此
\[
 \lambda_U(\Gamma)\ge
 \frac1{\sum_j\lambda_j^{-1}+\varepsilon}.
\]
令 \(\varepsilon\downarrow0\)，即得匹配的下界。
\end{proof}

并联相加的是倒极值长度。各通道的长度被归一化为同一个数 \(1\)，曲线可以选择其中任一条路线；面积却按不相交部分相加，这正是倒数出现的原因。串联比较则要求每条总曲线依次穿过所有子通道，长度贡献也必须相加。若实际允许曲线在不同通道之间穿梭，所研究的族就不再是上述并集，必须先重新辨认全部路线，不能仅凭图形看似并列套用公式。

\subsection{模量}
\label{b3:subsec:080403}

矩形中的计算使宽高比成为一个不依赖矩形映射选择的量。定义时仍要保留连接方向，否则一个数及其倒数都可能被叫作同一个四边形的模。

\Needspace{6\baselineskip}
\begin{definition}\label{b3:def:08-quadrilateral-modulus}
设 \(Q\subseteq\mathbb C\) 为有界 Jordan 区域，四个互异边界点 \(p_0,p_1,p_2,p_3\) 按正向排列。令 \(\Gamma_{\mathrm h}\) 为 \(Q\) 内局部可求长、两极限端点分别落在开弧 \((p_3,p_0)\) 与 \((p_1,p_2)\) 上的全部曲线族，开弧均沿正向边界选取。这两弧分别标记为左、右边。该标记四边形的\term[sibianxing-moliang]{四边形模量}[modulus of a quadrilateral]定义为
\[
 \operatorname{mod}(Q;p_0,p_1,p_2,p_3)
 \coloneqq\lambda_Q(\Gamma_{\mathrm h}).
\]
标记已经固定时简记为 \(\operatorname{mod}Q\)。
\end{definition}

\Needspace{6\baselineskip}
\begin{proposition}[矩形参数及标记唯一性]\label{b3:prop:08-marked-modulus-uniqueness}
设 \(Q\subseteq\mathbb C\) 为有界 Jordan 区域，\(p_0,p_1,p_2,p_3\in\partial Q\) 互异并按正向排列。令 \(\Gamma_{\mathrm h}\) 连接开弧 \((p_3,p_0)\) 与 \((p_1,p_2)\)，\(\Gamma_{\mathrm v}\) 连接开弧 \((p_0,p_1)\) 与 \((p_2,p_3)\)，两族均包括所有内部局部可求长且有这些极限端点的曲线。若 \(a,b>0\)，\(R_{a,b}=(0,a)\times(0,b)\)，双全纯映射 \(F:Q\rightarrow R_{a,b}\) 延为闭区域同胚，并将四个顶点依次映到 \(0,a,a+\mathrm{i}b,\mathrm{i}b\)，则
\[
 0<\operatorname{mod}Q=\frac ab<\infty,\quad
 \lambda_Q(\Gamma_{\mathrm v})=\frac ba
 =\frac1{\operatorname{mod}Q}.
\]
因此所有保持这些标记的矩形映射具有相同宽高比。目标矩形固定时，\(F\) 唯一；允许目标矩形等比例缩放时，映射也只差同一正比例因子。
\end{proposition}
\begin{proof}
边界同胚保证 \(F\) 将两对开边分别送到矩形的两对开边。\(F\) 及其逆在内部紧集上均为光滑映射，所以把局部可求长曲线互相对应；端点极限则由边界连续性保持。因此 \(F\Gamma_{\mathrm h}\)、\(F\Gamma_{\mathrm v}\) 恰为矩形中的全部水平、竖直方向连接族。共形不变性及\cref{b3:thm:08-rectangle-extremal}给出两个等式。若另一矩形映射的边长为 \(a',b'\)，同样的等式给 \(a'/b'=a/b\)，证明宽高比唯一。

设两映射 \(F_1,F_2\) 的目标矩形相同。\(F_2^{-1}\circ F_1\) 是 \(Q\) 的自同构，并在边界固定四个标记。用一个有边界同胚的 Riemann 映射将它共轭到圆盘，便得到固定三个以上边界点的圆盘自同构。圆盘自同构是 Möbius 变换，而非恒等 Möbius 变换至多有两个固定点，所以这个自同构为恒等，\(F_1=F_2\)。最后，若 \(a'/b'=a/b\)，取 \(c=a'/a=b'/b>0\)，则 \(c\cdot F_1\) 与 \(F_2\) 映到同一矩形并保持全部标记，由刚证明的唯一性可知二者相等。
\end{proof}

因而每个标记四边形都有唯一的归一化矩形 \((0,M)\times(0,1)\)，其中 \(M=\operatorname{mod}Q\)。两个标记 Jordan 四边形共形等价，当且仅当其模量相同：必要性来自共形不变性；充分性则是把两者映到同一个归一化矩形，再复合一个映射与另一个的逆。这里的等价要求逐一保持四个标记。如果把标记循环移一位，左右连接族变成原来的上下连接族，模量就变为 \(1/M\)；移两位只是倒转左右两端，曲线族不变。

\ProseParagraphBreak
\cref{b3:prop:08-marked-rectangle-existence}给出的具体参数满足
\[
 \operatorname{mod}Q=\frac{K(k)}{K(k')},\quad k'=\sqrt{1-k^2}.
\]
\cref{b3:prop:07-elliptic-aspect}说明这个比值随 \(k\in(0,1)\) 严格递增并遍历全部正数，所以模量也唯一确定归一化后的第四预顶点 \(k^{-2}\)。四点在圆周或实轴上的交比可以记录这一参数；原区域中四个欧氏坐标的交比却不能直接代替它，因为一般共形映射并不保持欧氏交比，只有 Möbius 变换保持。

\ProseParagraphBreak
矩形极值密度还给出原区域上的一个明确极值密度。若 \(F:Q\rightarrow R_{a,b}\) 保持标记，则将常密度拉回，得到 \(\rho(z)=|F'(z)|\)。共形换元给面积 \(A_\rho=ab\)，曲线长度对应给 \(L_\rho(\Gamma_{\mathrm h})=a\)，所以其商为 \(a/b\)。即使 \(F'\) 在边界附近无界，这仍是合法密度，因为总面积恰为有限的矩形面积。

\ProseParagraphBreak
一个可显式计算的例子是扇形通道
\[
 S=\{\,z:r<|z|<R,\ 0<\arg z<\theta\,\},\quad
 0<r<R,\quad0<\theta<2\pi.
\]
将两个圆弧标记为左、右边，两条径向边标记为下、上边。选定 \(0<\arg z<\theta\) 的分支后，\(F(z)=\log(z/r)\) 把它映到宽 \(\ell=\log(R/r)\)、高 \(\theta\) 的矩形，因而模量为 \(\ell/\theta\)。对应的极值密度 \(1/|z|\) 也能直接核算：其面积为
\[
 \int_0^\theta\int_r^R\frac1{s^2}\,s\,\mathrm ds\,\mathrm d\varphi
 =\theta\ell,
\]
任意两圆弧连接曲线的密度长度至少为 \(|\log R-\log r|=\ell\)，而径向段达到此值。因此其商为 \(\ell^2/(\theta\ell)=\ell/\theta\)。例如 \(r=1\)、\(R=e^3\)、\(\theta=\pi/2\) 时，左右连接模量为 \(6/\pi\)，另一个方向为 \(\pi/6\)。这里的左右边在原图上是内、外圆弧，进一步说明模量依赖指定的连接方向，而非边弧在纸面上的名字。

\subsection{边界估计}
\label{b3:subsec:080404}

精确矩形映射通常难以求出，但极值长度的定义允许只在必经的狭窄部分布置密度。密度以外的区域不耗费面积；只要每条允许曲线确实横穿这部分，局部长度就能给全族的下界。

\Needspace{6\baselineskip}
\begin{proposition}[必经矩形的密度下界]\label{b3:prop:08-bottleneck-density}
设 \(U\subseteq\mathbb C\) 为非空开集，\(\Gamma\) 为其中局部可求长曲线族。设 \(a,b>0\)、\(c,d\in\mathbb R\)，且
\[
 R=(c,c+a)\times(d,d+b)\subseteq U.
\]
假定每条 \(\gamma\in\Gamma\) 都包含一条子曲线，其内部位于 \(R\)，两极限端点分别在 \(R\) 的左、右开边上。则
\[
 \lambda_U(\Gamma)\ge\frac ab.
\]
\end{proposition}
\begin{proof}
在 \(R\) 上取密度 \(1\)，在 \(U\setminus R\) 上取 \(0\)，所得 Borel 密度的面积为 \(ab\)。每条 \(\gamma\) 的指定子曲线横坐标从 \(c\) 走到 \(c+a\)，其欧氏长度至少为 \(a\)。子曲线内部密度为 \(1\)，总曲线上其余部分的长度贡献非负，故 \(L_\rho(\Gamma)\ge a\)。将这个密度代入定义，即得 \(a^2/(ab)=a/b\) 的下界。
\end{proof}

这项假设要求每条路线横穿整个矩形。矩形只是包含在 \(U\) 中，或者某些曲线经过其中，都不能给同样的结论：曲线族的下确界还会考虑所有绕过它的路线。将局部几何转成边界估计时，先证明必经性，再选密度，次序不能颠倒。

\ProseParagraphBreak
变宽通道可以更细致地分辨瓶颈的贡献。对 \(\varepsilon>0\)，在 \(-1\le x\le1\) 上令
\[
 h_\varepsilon(x)=\frac{\varepsilon+|x|}{1+2\varepsilon},\quad
 Q_\varepsilon=\{\,(x,y):-1<x<1,\ 0<y<h_\varepsilon(x)\,\}.
\]
每个 \(Q_\varepsilon\) 都是 Jordan 区域，且面积完全相同：
\[
 |Q_\varepsilon|
 =\frac1{1+2\varepsilon}\int_{-1}^{1}(\varepsilon+|x|)\,\mathrm dx=1.
\]
中部的高度为 \(\varepsilon/(1+2\varepsilon)\)，两端则明显较宽。\cref{b3:prop:08-variable-channel}给左右连接族 \(\Gamma_\varepsilon\) 的下界
\[
 \begin{aligned}
 \lambda(\Gamma_\varepsilon)
 &\ge\int_{-1}^{1}\frac{\mathrm dx}{h_\varepsilon(x)}\\
 &=2(1+2\varepsilon)\int_0^1\frac{\mathrm dx}{\varepsilon+x}
 =2(1+2\varepsilon)\log\frac{1+\varepsilon}{\varepsilon}.
 \end{aligned}
\]
当 \(\varepsilon\downarrow0\) 时，右侧趋于无穷。因此，即使面积始终等于 \(1\)，模量也能任意大。宽为 \(2\)、高为 \(1/2\) 的矩形同样面积为 \(1\)，模量却恒为 \(4\)。例如 \(\varepsilon=1/4\) 时，瓶颈下界已经是 \(3\log5>4\)。我们无需求出 \(Q_\varepsilon\) 的矩形映射，就能严格区分这两个等面积通道。

\ProseParagraphBreak
用于这个下界的密度为 \(\rho_\varepsilon(x,y)=(1+2\varepsilon)/(\varepsilon+|x|)\)。它没有把全部面积预算集中到一条竖线，而是在一个实际宽度的区域内逐渐增大；积分同时计入窄度及其沿通道持续的距离。通道越接近中部夹断，两端之间的连接代价就越大。这样的密度下界是讨论边界退化时可直接使用的工具，单独的欧氏面积或最窄截面尺寸都不能提供同样完整的信息。

\subsection{拟共形观点}
\label{b3:subsec:080405}

共形变换在一点的所有方向上具有同一个伸长，长度平方与面积的变化因此精确抵消。一般光滑变换把小圆变成椭圆；长度可能取到长轴伸长，但面积仍由两根轴的乘积控制。如果两轴之比处处不超过 \(K\)，最后留下的畸变因子便是 \(K\)。

\Needspace{6\baselineskip}
\begin{proposition}[光滑拟共形映射的极值长度畸变]\label{b3:prop:08-smooth-quasiconformal-distortion}
设 \(U,V\subseteq\mathbb C\) 为区域，\(1\le K<\infty\)，\(\Gamma\) 为 \(U\) 内局部可求长曲线族。设 \(f:U\rightarrow V\) 为保定向 \(C^1\) 微分同胚。将实导数 \(Df(z)\) 的最大、最小奇异值记为 \(s_1(z)\ge s_2(z)>0\)，并假定对每个 \(z\in U\) 有 \(s_1(z)/s_2(z)\le K\)。则
\[
 K^{-1}\lambda_U(\Gamma)
 \le\lambda_V(f\Gamma)
 \le K\lambda_U(\Gamma).
\]
这里 \(f\Gamma=\{\,f\circ\gamma:\gamma\in\Gamma\,\}\)，不等式也允许极值长度为 \(0\) 或 \(+\infty\)。
\end{proposition}
\begin{proof}
实矩阵 \(Df(z)^{\mathsf T}Df(z)\) 对称正定，其两个特征值为 \(s_1(z)^2,s_2(z)^2\)。故对任意实向量 \(v\) 有
\[
 s_2(z)\cdot|v|\le|Df(z)v|\le s_1(z)\cdot|v|.
\]
又因 \(f\) 保定向，其 Jacobi 行列式满足
\[
 J_f(z)=\det Df(z)=s_1(z)\cdot s_2(z)>0,\quad
 J_f(z)\le s_1(z)^2\le K\cdot J_f(z).
\]
后一关系把最大方向的长度伸长与真实面积伸长联系起来。

给定 \(V\) 上的非负 Borel 密度 \(\sigma\)，假设 \(0<A_\sigma<\infty\)。在 \(U\) 上定义
\[
 \rho(z)=\sigma(f(z))\cdot s_1(z).
\]
对于 \(\gamma\) 的一条紧子弧，其像是 \(U\) 内的紧集，\(f\) 在它的一个相对紧邻域上导数有界，所以复合曲线仍可求长。使用该子弧的弧长参数以及绝对连续链式法则，得到
\[
 \int_{f\circ\gamma}\sigma\,\mathrm ds
 =\int_\gamma\sigma(f(z))\cdot|Df(z)\tau(z)|\,\mathrm ds
 \le\int_\gamma\sigma(f(z))\cdot s_1(z)\,\mathrm ds,
\]
其中 \(\tau\) 是几乎处处存在的单位切向量。先在紧子弧上写此式，再以单调极限穷竭整条曲线，便得到 \(L_\sigma(f\circ\gamma)\le L_\rho(\gamma)\)，包括长度为无穷的情况。对 \(\gamma\in\Gamma\) 取下确界，有
\[
 L_\sigma(f\Gamma)\le L_\rho(\Gamma).
\]

平面换元公式及奇异值关系给出
\[
 \begin{aligned}
 A_\sigma
 &=\int_U\sigma(f(z))^2\cdot J_f(z)\,\mathrm dA(z)\\
 &\le A_\rho
 =\int_U\sigma(f(z))^2\cdot s_1(z)^2\,\mathrm dA(z)\\
 &\le K\int_U\sigma(f(z))^2\cdot J_f(z)\,\mathrm dA(z)
 =K A_\sigma.
 \end{aligned}
\]
因而 \(\rho\) 具有有限正面积，确实是源区域的合法竞争密度。如果 \(\lambda_U(\Gamma)=\infty\)，所需上界自动成立；否则 \(L_\rho(\Gamma)\) 有限，并可计算
\[
 \frac{L_\sigma(f\Gamma)^2}{A_\sigma}
 \le\frac{L_\rho(\Gamma)^2}{A_\sigma}
 \le K\frac{L_\rho(\Gamma)^2}{A_\rho}
 \le K\lambda_U(\Gamma).
\]
对 \(\sigma\) 取上确界，得到右侧不等式。

在 \(w=f(z)\) 处，\(D(f^{-1})(w)=Df(z)^{-1}\) 的两个奇异值为 \(1/s_2(z)\ge1/s_1(z)\)，其比仍不超过 \(K\)。逆映射同样保定向且为 \(C^1\)；它把 \(f\Gamma\) 恢复为 \(\Gamma\)。因此将已证上界用于 \(f^{-1}\)，得到 \(\lambda_U(\Gamma)\le K\lambda_V(f\Gamma)\)，也就是左侧不等式。
\end{proof}

面积中出现 \(s_1s_2\)，长度平方中出现 \(s_1^2\)，两者相除留下 \(s_1/s_2\)，所以常数是 \(K\) 而非 \(K^2\)。当 \(K=1\) 时，两方向伸长相同，结论恢复共形不变性。若还已知 \(f\) 在两个 Jordan 区域的边界间给出保持四个标记的同胚，那么它把全部指定连接曲线族互相对应，以上不等式就直接成为四边形模量的双边比较。

\ProseParagraphBreak
这个证明使用了 \(C^1\) 微分同胚的线积分链式法则与面积换元。一般拟共形同胚只具有相应的 Sobolev 正则性，导数及伸长条件通常几乎处处成立，还需要处理曲线族的零模例外。相关推广将在\appref{b3:app:B}中说明；不能仅把本命题的逐点假设改成几乎处处，就视为已经完成一般情形的证明。

\ProseParagraphBreak
径向伸缩给出取等模型，而且能同时看到两种曲线方向的相反变化。对 \(\alpha>0\)，令
\[
 F_\alpha(z)=z|z|^{\alpha-1},\quad z\ne0.
\]
它在极坐标中将 \((s,\theta)\) 变成 \((s^\alpha,\theta)\)，逆为 \(F_{1/\alpha}\)。径向、角向单位向量分别被伸长 \(\alpha s^{\alpha-1}\)、\(s^{\alpha-1}\)，且像仍正交，因此
\[
 J_{F_\alpha}=\alpha s^{2\alpha-2}>0,\quad
 K=\max\{\alpha,\alpha^{-1}\}.
\]
这证明 \(F_\alpha\) 是去心平面的光滑保向微分同胚，并给出它的确切伸长比。对 \(0<r<R<\infty\)，有
\[
 F_\alpha(A(r,R))=A(r^\alpha,R^\alpha),\quad
 \operatorname{mod}F_\alpha(A(r,R))
 =\frac{\log(R^\alpha/r^\alpha)}{2\pi}
 =\alpha\operatorname{mod}A(r,R).
\]
模量是两端连接族的极值长度，所以 \(\alpha\ge1\) 时达到上界，\(0<\alpha\le1\) 时达到下界。绕行族的极值长度则乘以 \(1/\alpha\)，达到另一个端点。由此可见，同一个非共形映射对不同方向的曲线族可以产生相反的变化，畸变因子不能整体改善。角坐标保持不变也不表示所有夹角保持不变：除 \(\alpha=1\) 外，径向与角向伸长不同，微分就不是一个复数的乘法。


\begin{chaptersummary}
极值长度把一项几何问题分成两个方向：选择一个密度，证明每条允许曲线都有足够长度，便得到下界；选择一族便于平均的允许曲线，证明任意密度中总有一条不太长，便得到上界。圆环中的 \(1/|z|\) 与矩形中的常密度使这两步相遇，分别给出 \(\log(R/r)/(2\pi)\) 与指定方向的宽高比。圆环的绕行族、四边形的另一对边连接族给出相应倒数；换了曲线族，不能只沿用原来的数值名称。

\ProseParagraphBreak

比较时须同时保留区域与路径条件。本质包含保持环域的孔，串联要求每条路径逐段穿过，并联要求路径在不交通道中择一而行。试探势把所有路径的长度下界归结为两端势差，其能量给出可优化的竞争值。光滑拟共形变换中，长度按导数伸长变化，面积按两个主伸长的乘积变化，因而留下一个 \(K\) 的畸变因子。面积几乎处处相等的密度未必对全部曲线等效；无限模量也未必唯一决定退化区域。精确计算、比较不等式与分类的存在输入各自解决一部分问题，不能互相替代。
\end{chaptersummary}

\BookExercises

% 本章8题均为自拟改写；训练几何辨认、畸变最优化、密度代表、通道分配与试探势。
\Needspace{6\baselineskip}
\begin{exercise}\label{b3:ex:08-1}
设 \(R>1\)、\(a,b\in\mathbb C\)、\(a\ne0\)。判断
\[
 U=\{\,z\in\mathbb C\mid1<|az+b|<R\,\}
\]
是否为真正的椭圆环，求其中心、边界半径和模量。再设 \(\alpha,\beta>0\)，令
\[
 T(x+\mathrm{i}y)=\alpha x+\mathrm{i}\beta y,\quad
 E=T(A(1,R)).
\]
写出 \(E\) 的边界方程，判断 \(T\) 何时共形，并利用最大、最小伸长给 \(\operatorname{mod}E\) 一个双边界。能否仅凭这两项伸长就把矩形的宽高比变换公式当作 \(E\) 的精确模量公式？
\end{exercise}
\begin{solution}
将复仿射式配成中心形式，有
\[
 |az+b|=|a|\cdot|z+b/a|.
\]
所以 \(U\) 是以 \(-b/a\) 为中心、内外半径分别为 \(1/|a|\)、\(R/|a|\) 的同心圆环。映射 \(z\mapsto az+b\) 的导数为非零复常数，将它双全纯映到 \(A(1,R)\)，故
\[
 \operatorname{mod}U=\frac{\log R}{2\pi}.
\]
复数乘法在实平面上只有旋转与等向缩放，两条主轴不能独立伸缩。

写 \(w=X+\mathrm{i}Y=T(z)\)，则
\[
 E=\left\{\,X+\mathrm{i}Y\,\middle|\,
 1<\frac{X^2}{\alpha^2}+\frac{Y^2}{\beta^2}<R^2\,\right\}.
\]
两条边界的半轴分别是 \((\alpha,\beta)\) 与 \((R\alpha,R\beta)\)。当 \(\alpha\ne\beta\) 时它们才是真正的椭圆。实导数为对角矩阵 \(\operatorname{diag}(\alpha,\beta)\)，行列式 \(\alpha\beta>0\)，两项主伸长就是 \(\alpha,\beta\)。它是复线性映射当且仅当 \(\alpha=\beta\)，亦可由 Cauchy--Riemann 条件直接核对。

置
\[
 \kappa=\max\left\{\frac\alpha\beta,\frac\beta\alpha\right\}.
\]
映射 \(T\) 及其逆均光滑，并将两条边界分别对应，所以圆环的全部两端连接曲线恰好送到椭圆环的连接族。由光滑拟共形畸变及圆环计算，
\[
 \frac1\kappa\frac{\log R}{2\pi}
 \le\operatorname{mod}E
 \le\kappa\frac{\log R}{2\pi}.
\]
这里给出的是界，未辨认椭圆环的极值密度。矩形在对角伸缩后仍是矩形，新的宽高比可直接计算；同一论证不能把椭圆环也当成标准圆环。尤其同一个椭圆环旋转90度会交换 \(\alpha,\beta\)，其环域模量不变；若无条件写成 \((\alpha/\beta)\log R/(2\pi)\)，交换两轴后就会给出不同答案。
\end{solution}

\Needspace{6\baselineskip}
\begin{exercise}\label{b3:ex:08-2}
设 \(0<r<R<\infty\)、\(0<s<S<\infty\)。在所有保向 \(C^1\) 微分同胚
\[
 f:A(r,R)\longrightarrow A(s,S)
\]
中，要求 \(f\) 连续延为闭圆环的同胚，并分别将内、外圆周对应。记其最大伸长比为
\[
 K(f)=\sup_{z\in A(r,R)}\frac{s_1(Df(z))}{s_2(Df(z))}.
\]
求 \(K(f)\) 的最小可能值，并构造达到该值的映射。如果也允许交换内外圆周，最小值是否改变？
\end{exercise}
\begin{solution}
记两个对数宽度及其比为
\[
 L=\log(R/r),\quad M=\log(S/s),\quad\eta=\frac ML>0.
\]
边界同胚与内部微分同胚使两端连接曲线族一一对应。若 \(K(f)<\infty\)，由\cref{b3:prop:08-smooth-quasiconformal-distortion}，
\[
 K(f)^{-1}\frac{L}{2\pi}
 \le\frac{M}{2\pi}
 \le K(f)\frac{L}{2\pi}.
\]
所以 \(K(f)\ge\max\{\eta,\eta^{-1}\}\)。伸长比无界的映射也不会降低这个下界。

反向构造取
\[
 f_0(z)=s\left(\frac{|z|}{r}\right)^\eta\frac z{|z|}.
\]
极坐标中，角度保持不变，半径 \(t\) 被送到
\(s(t/r)^\eta\)。它在 \([r,R]\) 上严格递增，并将两端送到 \(s,S\)，因为 \((R/r)^\eta=S/s\)。逆半径函数为 \(r(u/s)^{1/\eta}\)，所以 \(f_0\) 具有题设要求的光滑逆与边界延拓。

在源点半径 \(t\) 处，径向单位向量与角向单位向量分别被伸长
\[
 \frac{s}{r^\eta}\eta t^{\eta-1},\quad
 \frac{s}{r^\eta}t^{\eta-1},
\]
两像仍正交。行列式为正，伸长比处处等于 \(\max\{\eta,\eta^{-1}\}\)。故最小值不仅有下界，而且由 \(f_0\) 达到：
\[
 \min_f K(f)=
 \max\left\{\frac{\log(S/s)}{\log(R/r)},
             \frac{\log(R/r)}{\log(S/s)}\right\}.
\]
例如从 \(A(1,e^2)\) 到 \(A(1,e^3)\)，最小值为 \(3/2\)，由 \(z\mapsto z|z|^{1/2}\) 达到。

目标圆环的倒数映射 \(J(w)=sS/w\) 双全纯且交换两端，其导数在圆环内非零。对一个交换两端的 \(f\)，\(J\circ f\) 保持两端，且共形的后复合不改变伸长比。反过来也可把达到最小值的 \(f_0\) 与 \(J\) 复合。因此允许交换两端时最小值不变。
\end{solution}

\Needspace{6\baselineskip}
\begin{exercise}\label{b3:ex:08-3}
在 \(A=A(1,e)\) 中，令 \(\Gamma_{\mathrm r}\) 为全部两端连接曲线，\(\Gamma_{\mathrm c}\) 为绕原点指数为 \(1\) 或 \(-1\) 的全部可求长闭曲线。记
\[
 E_{\mathrm r}=\{\,x\in\mathbb R\mid1<x<e\,\},\quad
 E_{\mathrm c}=\{\,z\mid|z|=\sqrt e\,\},\quad
 \rho(z)=\frac1{|z|}.
\]
分别把 \(\rho\) 在 \(E_{\mathrm r}\) 或 \(E_{\mathrm c}\) 上改为零，得到密度 \(\rho_{\mathrm r}\)、\(\rho_{\mathrm c}\)。求它们的面积、对两族的最短长度及对应长度面积商；再判断同时在两集合上改为零时会怎样。这些计算是否改变两曲线族本身的极值长度？
\end{exercise}
\begin{solution}
两集合都为零面积 Borel 集，故三种修改密度的面积仍是
\[
 \int_A\rho^2\,\mathrm{d}A
 =2\pi\int_1^e\frac{\mathrm{d}t}{t}=2\pi.
\]
\(E_{\mathrm r}\) 本身是一条径向连接曲线，其 \(\rho_{\mathrm r}\) 长度为零。因此
\(L_{\rho_{\mathrm r}}(\Gamma_{\mathrm r})=0\)。类似地，圆周 \(E_{\mathrm c}\) 在 \(\rho_{\mathrm c}\) 下长为零，所以
\(L_{\rho_{\mathrm c}}(\Gamma_{\mathrm c})=0\)。

另一个方向需核对每条曲线，不能只计算几个代表。先用一个实分析事实：若绝对连续实函数 \(h\) 在集合 \(E\) 上等于同一个常数，则 \(h'=0\) 在 \(E\) 上几乎处处成立。在 \(E\) 中非孤立的可微点，用 \(E\) 内的点列趋近可使差商为零；孤立点至多可数，而不可微点集为零测集。

对一条绕行曲线作弧长参数化 \(\gamma\)，取其连续辐角 \(\theta\)。曲线像是 \(A\) 内的紧集，沿有限个局部对数分支使用链式法则，得到 \(\theta\) 绝对连续，且几乎处处满足
\[
 |\theta'|\le\frac{|\gamma'|}{|\gamma|}.
\]
当 \(\gamma\) 落在正实轴线段 \(E_{\mathrm r}\) 时，\(\theta\) 属于可数集合 \(2\pi\mathbb Z\)。对每个水平集应用上述事实，可知这些参数处 \(\theta'=0\) 几乎处处。删去该线段上的密度因而不会删去辐角变化，仍有
\[
 L_{\rho_{\mathrm r}}(\gamma)
 \ge\int|\theta'|\ge|\theta(\text{终点})-\theta(\text{起点})|
 =2\pi.
\]
任意同心圆周只在一个点遇到删去的径线，长度仍为 \(2\pi\)，故
\(L_{\rho_{\mathrm r}}(\Gamma_{\mathrm c})=2\pi\)。

对任一两端连接曲线，在每条紧子弧上考察
\(h=\log|\gamma|\)。它绝对连续，且
\(|h'|\le|\gamma'|/|\gamma|\)。曲线落在 \(E_{\mathrm c}\) 时 \(h=1/2\)，故这些参数上的 \(h'\) 几乎处处为零。让紧子弧的两端趋向两个圆周，得到
\[
 L_{\rho_{\mathrm c}}(\gamma)\ge|1-0|=1.
\]
径向直线只在一个点遇到删去的圆周，其长度为1，所以
\(L_{\rho_{\mathrm c}}(\Gamma_{\mathrm r})=1\)。

于是两种固定密度给出的资料是
\[
\begin{array}{c|cc|cc}
 &L(\Gamma_{\mathrm r})&L(\Gamma_{\mathrm c})
 &L(\Gamma_{\mathrm r})^2/A&L(\Gamma_{\mathrm c})^2/A\\ \hline
 \rho_{\mathrm r}&0&2\pi&0&2\pi\\
 \rho_{\mathrm c}&1&0&1/(2\pi)&0
\end{array}
\]
若两集合上同时置零，两族都含一条零长度曲线，所以两个商都为零，面积仍为 \(2\pi\)。这些都是特定密度的竞争值；区域与曲线族没有改变，极值长度仍分别为 \(1/(2\pi)\) 和 \(2\pi\)。原密度 \(1/|z|\) 仍是允许的竞争者，不能把某个代表失去极值性误说成曲线族的不变量改变。
\end{solution}

\Needspace{6\baselineskip}
\begin{exercise}\label{b3:ex:08-4}
设两条不相交的水平矩形通道 \(Q_1,Q_2\) 分别宽6、高2及宽4、高1，\(\Gamma_j\) 是各自的全部左右连接曲线。曲线可任选一条通道，即总族为 \(\Gamma_1\cup\Gamma_2\)。在总密度面积归一化为1时，求最大的最短长度、相应的两个通道面积分配，并构造达到该值的密度。

\ProseParagraphBreak

另设两通道位于某个开集内，规定每条总路径都必须包含分别横穿 \(Q_1,Q_2\) 的子曲线。在两个矩形各取常密度、其余处取零，仍以总面积1为预算。怎样分配可使由这两段横穿得到的长度下界最大？比较两种分配，并说明串联估计是否必为整个区域的精确极值长度。
\end{exercise}
\begin{solution}
矩形计算给两个单独曲线族的极值长度为3和4。对总面积为1的密度，记两通道面积为 \(A_1,A_2\)，最短长度为 \(L_1,L_2\)。若 \(A_j>0\)，定义给
\[
 L_1^2\le3A_1,\quad L_2^2\le4A_2.
\]
若某个 \(A_j=0\)，沿该矩形的水平截面使用 Tonelli 定理，几乎所有截面上的平方密度积分为零，因而有水平连接线长为零。故对应 \(L_j=0\)，上述不等式仍有效，不能只凭面积为零就断言每条路径长度都为零。

并联时最短长度是 \(L=\min\{L_1,L_2\}\)，故
\[
 L^2\le\min\{3A_1,4A_2\},\quad A_1+A_2=1.
\]
当 \(3A_1=4A_2\) 时右端最大，否则把少量面积从较大项转向较小项就会改进最小值。因此最优分配和最优值为
\[
 A_1=\frac47,\quad A_2=\frac37,\quad L^2\le\frac{12}{7}.
\]
在 \(Q_1,Q_2\) 上分别取常密度
\[
 c_1=\frac1{\sqrt{21}},\quad c_2=\frac3{2\sqrt{21}}.
\]
面积为 \(12c_1^2+4c_2^2=1\)，各族的最短长度分别为 \(6c_1,4c_2\)，且二者都等于 \(6/\sqrt{21}=\sqrt{12/7}\)。水平线达到这些长度，而任何左右连接路径的欧氏长度至少为宽度，所以构造对全部曲线有效。总族的极值长度为 \(12/7\)。

串联时，非负常密度 \(d_1,d_2\) 给每条总路径的长度下界 \(6d_1+4d_2\)，面积为 \(12d_1^2+4d_2^2\)。Cauchy--Schwarz 给
\[
 (6d_1+4d_2)^2
 \le(12d_1^2+4d_2^2)
 \left(\frac{6^2}{12}+\frac{4^2}{4}\right)
 =7(12d_1^2+4d_2^2).
\]
等号要求 \(d_1:d_2=1:2\)。在面积1的约束下，取
\[
 d_1=\frac1{2\sqrt7},\quad d_2=\frac1{\sqrt7},
\]
便得到最强的这类保证 \(L\ge\sqrt7\)，因此总族极值长度至少为7。两个矩形占用的面积现在分别为 \(3/7,4/7\)，与并联的 \(4/7,3/7\) 相反。

并联配置要使两条备选路径一样难走；串联配置则使必经段的长度贡献相加。串联条件只给总路径中一部分的约束，连接两矩形的其余区域可能增加阻隔，故仅据题设不能断言整个曲线族的极值长度等于7。
\end{solution}

\Needspace{6\baselineskip}
\begin{exercise}\label{b3:ex:08-5}
设
\[
 Q=\{\,x+\mathrm{i}y\mid0<x<1,\ 0<y<1+x\,\},
\]
\(\Gamma\) 为连接左右两条竖直边的全部曲线。先用变宽通道与内含矩形给 \(\lambda(\Gamma)\) 的上下界，再利用直线段族
\[
 \gamma_v(x)=x+\mathrm{i}v(1+x),\quad
 0<x<1,\quad0<v<1
\]
改进上界，证明
\[
 \log2\le\lambda(\Gamma)\le\frac{4\log2}{\pi}<1.
\]
最后比较它与同宽、同面积矩形的左右连接模量，说明仅有面积不能确定这项共形参数。
\end{exercise}
\begin{solution}
高度函数 \(h(x)=1+x\) 连续，且 \(1\le h\le2\)。\Cref{b3:prop:08-variable-channel}给出
\[
 \lambda(\Gamma)\ge\int_0^1\frac{\mathrm{d}x}{1+x}=\log2.
\]
单位正方形 \(Q_0=(0,1)\times(0,1)\) 包含于 \(Q\)，其左右连接曲线也是 \(\Gamma\) 中的曲线。使用\cref{b3:prop:08-ambient-independence}把这个子族放到同一环境，再用曲线族包含原则，得到 \(\lambda(\Gamma)\le1\)。

为改进上界，注意 \(\gamma_v\) 在内部始终满足 \(0<v(1+x)<1+x\)，左端是 \((0,v)\)，右端是 \((1,2v)\)，确实连接题目指定的两条竖边。映射
\[
 (x,v)\longmapsto(x,(1+x)v)
\]
将 \((0,1)^2\) 一一映到整个 \(Q\)，Jacobi 行列式为 \(1+x>0\)。因此这一参数族既是允许的曲线族，又给出了可用于面积平均的全域坐标。

任取有限正面积 Borel 密度 \(\sigma\)，记总面积为 \(A\)、全族最短长度为 \(L\)。沿 \(\gamma_v\) 的速度为 \(\sqrt{1+v^2}\)，故其长度 \(\ell(v)\) 满足
\[
 \begin{aligned}
 \ell(v)^2
 &=(1+v^2)\left(\int_0^1
       \sigma(x,(1+x)v)\,\mathrm{d}x\right)^2\\
 &\le(1+v^2)\log2\int_0^1
       \sigma(x,(1+x)v)^2(1+x)\,\mathrm{d}x.
 \end{aligned}
\]
这里在 Cauchy--Schwarz 中配对 \(\sigma\sqrt{1+x}\) 与 \((1+x)^{-1/2}\)。换元与 Tonelli 定理保证右端的内积分对几乎所有 \(v\) 有限，且对 \(v\) 积分等于 \(A\)。每条 \(\gamma_v\) 都属于 \(\Gamma\)，所以 \(L\le\ell(v)\)。将上式除以 \(1+v^2\) 并积分，得到
\[
 L^2\int_0^1\frac{\mathrm{d}v}{1+v^2}
 \le A\log2,
 \quad\text{即}\quad
 \frac{L^2}{A}\le\frac{4\log2}{\pi}.
\]
取密度上确界得到改进后的上界。对 \(0<x<1\)，\(1+x^2<1+x\)，故
\(\pi/4=\int_0^1(1+x^2)^{-1}\,\mathrm{d}x>\log2\)，上界确实小于1。

通道面积为 \(\int_0^1(1+x)\,\mathrm{d}x=3/2\)。同宽、同面积矩形的高度为 \(3/2\)，左右连接模量为 \(2/3\)。严格的 Cauchy--Schwarz 不等式给
\[
 1<\left(\int_0^1(1+x)\,\mathrm{d}x\right)
       \left(\int_0^1\frac{\mathrm{d}x}{1+x}\right)
 =\frac32\log2,
\]
因为 \(1+x\) 不是常数。因此 \(\lambda(\Gamma)\ge\log2>2/3\)，两者面积相同而模量不同。上述两个界也没有相等，故尚未计算出 \(Q\) 的精确模量。
\end{solution}

\Needspace{6\baselineskip}
\begin{exercise}\label{b3:ex:08-6}
在去心盘 \(\mathbb D^*\) 中，令 \(\Gamma_0\) 为所有局部可求长、一个极限端点为0、另一个极限端点在单位圆周上的曲线。构造一个处处有限的非负 Borel 密度，使它的面积有限且为正，而每条 \(\Gamma_0\) 中曲线的长度都为无穷，从而直接证明 \(\lambda(\Gamma_0)=\infty\)。另用可去奇点与最大模原理证明：\(\mathbb D^*\) 不与任何有限圆环 \(A(r,R)\)、\(0<r<R<\infty\)，双全纯等价。
\end{exercise}
\begin{solution}
取
\[
 \rho(z)=\frac1{|z|\log(e/|z|)},\quad0<|z|<1.
\]
它在去心盘内处处为有限正数。极坐标换元，再取 \(t=\log(e/s)\)，得到
\[
 A_\rho=2\pi\int_0^1
 \frac{\mathrm{d}s}{s\log(e/s)^2}
 =2\pi\int_1^\infty\frac{\mathrm{d}t}{t^2}=2\pi.
\]
定义径向函数
\[
 u(z)=\log\log(e/|z|).
\]
它在去心盘光滑，满足 \(|\nabla u|=\rho\)；靠近穿孔时 \(u\to\infty\)，靠近单位圆周时 \(u\to0\)。对任一连接曲线的紧子弧，用弧长参数与链式法则可得
\[
 \int_{\text{子弧}}\rho\,\mathrm{d}s
 \ge|u(\text{一个端点})-u(\text{另一个端点})|.
\]
让两端分别趋向题设的极限端点，右侧可任意大，所以每条连接曲线都具有无穷 \(\rho\) 长度。这一个允许的有限正面积密度已经给出无穷的长度面积商，故 \(\lambda(\Gamma_0)=\infty\)。密度本身没有在任何内部点取无穷；无限长度来自靠近穿孔时持续累积的径向积分。

再假设有双全纯映射 \(h:\mathbb D^*\to A(r,R)\)。由 \(|h|<R\) 及\cref{b3:thm:03-removable}，\(h\) 在0可去，记延拓值为 \(a\)。因为 \(|h|>r\)，极限满足 \(r\le|a|\le R\)，特别地 \(a\ne0\)。若 \(|a|=R\)，延拓后的 \(h\) 在内部取得最大模，只能为常数；若 \(|a|=r\)，对无零的 \(h\) 应用 \(1/h\) 的最大模原理，仍得到常数。两者都与双全纯矛盾，故 \(a\in A(r,R)\)。

设逆映射为 \(g:A(r,R)\to\mathbb D^*\)。由 \(g(h(z))=z\)，让 \(z\to0\)，连续性给出 \(g(a)=0\)，违反逆映射的值域为去心盘。这个论证区分的是有限圆环与点状内端，不需要假定任意双全纯映射已经延为普通边界同胚。
\end{solution}

\Needspace{6\baselineskip}
\begin{exercise}\label{b3:ex:08-7}
设 \(Q=(0,a)\times(0,b)\)，其中 \(a,b>0\)。令 \(D\subset(0,b)\) 为可数稠密集，\(\Gamma_D\) 只由高度 \(y\in D\) 的完整水平线段组成；\(\Gamma_{\mathrm h}\) 则是矩形的全部左右连接曲线。证明
\[
 \lambda_Q(\Gamma_D)=\infty,\quad
 \lambda_Q(\Gamma_{\mathrm h})=a/b.
\]
第一式要求用一列有界 Borel 密度证明。若在 \(\Gamma_D\) 的极值问题中，把竞争密度限制为连续到闭矩形的非负函数，证明所得上确界反而为 \(a/b\)。解释为什么路径的稠密性与密度的 \(L^2\) 逼近都不能消除这个差别。
\end{exercise}
\begin{solution}
设
\[
 N=\bigcup_{y\in D}(0,a)\times\{y\}.
\]
这是零面积 Borel 集，并且在 \(Q\) 中稠密。对整数 \(n\ge1\)，定义有界 Borel 密度
\[
 \rho_n(z)=\begin{cases}n,&z\in N,\\1,&z\notin N.\end{cases}
\]
其面积始终为 \(ab\)。但 \(\Gamma_D\) 的每条曲线都完全落在 \(N\)，长度都是 \(na\)，因此
\[
 \frac{L_{\rho_n}(\Gamma_D)^2}{A_{\rho_n}}
 =n^2\frac ab\longrightarrow\infty.
\]
全部左右连接族的极值长度则由正文的矩形计算等于 \(a/b\)。两个族的包含方向是
\(\Gamma_D\subset\Gamma_{\mathrm h}\)，故无穷值与有限值正符合较小曲线族具有较大极值长度的原则。

现在取 \(\rho\in C(\overline Q)\)、\(\rho\ge0\)，且面积 \(A_\rho>0\)。水平线长度函数
\[
 h(y)=\int_0^a\rho(x,y)\,\mathrm{d}x,\quad0\le y\le b,
\]
由闭矩形上的一致连续性而连续。由于 \(D\) 在 \((0,b)\) 稠密，它在 \([0,b]\) 的闭包也是整个区间，故
\[
 L_\rho(\Gamma_D)=\inf_{y\in D}h(y)
 =\min_{0\le y\le b}h(y).
\]
这里连续到闭矩形的条件使边界高度也可通过内部高度逼近；不能对任意 Borel 密度沿用此等式。对每个 \(y\) 作 Cauchy--Schwarz，再对高度积分，得到
\[
 bL_\rho(\Gamma_D)^2
 \le\int_0^b h(y)^2\,\mathrm{d}y
 \le a\int_0^b\int_0^a\rho(x,y)^2\,\mathrm{d}x\,\mathrm{d}y
 =aA_\rho.
\]
所以限制后的上确界至多为 \(a/b\)，常密度1又达到此值。

\(\rho_n\) 与常密度1在二维面积几乎处处相等，因而它们作为 \(L^2\) 元素甚至完全相同，指定路径上的长度却不同。这里的密度代表没有足够的连续性，使稠密高度上的路径积分能够控制其他高度。将竞争类改成连续函数，改变了这个曲线族的极值问题，不能只以连续函数在 \(L^2\) 中稠密为由宣布上确界不变。
\end{solution}

\Needspace{6\baselineskip}
\begin{exercise}\label{b3:ex:08-8}
将单位正方形 \(Q=(0,1)^2\) 作线性剪切
\[
 T(x+\mathrm{i}y)=x+y+\mathrm{i}y,
\]
得到平行四边形
\[
 P=\{\,X+\mathrm{i}Y\mid0<Y<1,\ Y<X<Y+1\,\}.
\]
以 \(X=Y\) 与 \(X=Y+1\) 为两条指定连接边，记相应模量为 \(m\)。先计算 \(T\) 的伸长比，给出拟共形粗界，再用全部水平线段证明 \(m\le1\)。对实参数 \(c\)，考虑试探势
\[
 u_c(X,Y)=X-Y+c\sin\bigl(\pi(X-Y)\bigr)(Y-1/2).
\]
核对其边界值，计算能量并优化 \(c\)，证明
\[
 \frac1{\displaystyle2-\frac{48}{\pi^2(\pi^2+6)}}
 \le m\le1.
\]
比较这个下界与直接取线性势 \(u_0=X-Y\) 的下界，并说明是否已经求出精确模量。
\end{exercise}
\begin{solution}
实导数及其平方伸长矩阵为
\[
 DT=\begin{pmatrix}1&1\\0&1\end{pmatrix},\quad
 (DT)^{\mathsf T}DT=\begin{pmatrix}1&1\\1&2\end{pmatrix}.
\]
后者的特征值为 \((3\pm\sqrt5)/2\)，乘积为1。最大、最小伸长是它们的平方根，故伸长比为
\[
 K=\frac{3+\sqrt5}{2}.
\]
剪切保向且把指定两边对应，正方形模量为1，所以光滑拟共形界给
\(K^{-1}\le m\le K\)。

水平线段 \(Y=y\)、\(y<X<y+1\) 的每一个端点都落在题目指定的边上，它们构成允许的连接子族。给定 \(P\) 上有限正面积密度 \(\sigma\)，记其面积为 \(A\)、全族最短长度为 \(L\)。对几乎所有 \(y\)，水平段长 \(h(y)\) 满足
\[
 L^2\le h(y)^2
 =\left(\int_y^{y+1}\sigma(X,y)\,\mathrm{d}X\right)^2
 \le\int_y^{y+1}\sigma(X,y)^2\,\mathrm{d}X.
\]
对 \(0<y<1\) 积分得到 \(L^2\le A\)，所以 \(m\le1\)。这一上界已经强于粗界 \(K\)。

试探势 \(u_c\) 在整个实平面光滑。左边 \(X-Y=0\) 时为0，右边 \(X-Y=1\) 时为1；因此只需计算其能量，就可应用\cref{b3:prop:08-test-potential}。改用正方形坐标 \(X=x+y,Y=y\)，令
\[
 v_c(x,y)=u_c(x+y,y)=x+c\sin(\pi x)(y-1/2).
\]
换元的行列式为1，并且
\[
 (u_c)_X=(v_c)_x,\quad
 (u_c)_Y=(v_c)_y-(v_c)_x.
\]
所以物理坐标中的能量为
\[
 E(c)=\int_0^1\int_0^1
 \left(2(v_c)_x^2-2(v_c)_x(v_c)_y+(v_c)_y^2\right)
 \,\mathrm{d}x\,\mathrm{d}y.
\]
此处必须保留剪切引入的交叉项，不能把能量误作 \(\int_Q((v_c)_x^2+(v_c)_y^2)\)。

具体地，
\[
 (v_c)_x=1+c\pi\cos(\pi x)(y-1/2),\quad
 (v_c)_y=c\sin(\pi x).
\]
利用 \(\int_0^1(y-1/2)\,\mathrm{d}y=0\)、\(\int_0^1(y-1/2)^2\,\mathrm{d}y=1/12\)，以及正弦、余弦的基本积分，得到
\[
 \int_Q(v_c)_x^2=1+\frac{\pi^2c^2}{24},\quad
 \int_Q(v_c)_x(v_c)_y=\frac{2c}{\pi},\quad
 \int_Q(v_c)_y^2=\frac{c^2}{2}.
\]
因此
\[
 E(c)=2-\frac{4c}{\pi}+\frac{\pi^2+6}{12}c^2.
\]
这是一元严格凸二次式，其最小点与最小值为
\[
 c_* =\frac{24}{\pi(\pi^2+6)},\quad
 E(c_*)=2-\frac{48}{\pi^2(\pi^2+6)}.
\]
能量有限且为正；例如沿每条水平段，\(u_c\) 的端值差为1，Cauchy--Schwarz 已给 \(\int_P|(u_c)_X|^2\ge1\)，从而 \(E(c)\ge1\)。试探势命题于是给 \(m\ge1/E(c_*)\)，与前面的上界合成题设双边界。

线性势对应 \(c=0\)、\(E(0)=2\)，只能给 \(m\ge1/2\)。因 \(E(c_*)<2\)，新下界严格改进它，也强于 \(1/K\)。优化只在所给的一参数势族中进行，尚未证明这个势在所有允许函数中能量最小；下界与上界也未相等。因此所得是改进的模量估计，不是精确模量公式。
\end{solution}
